% GENERATED FILE. Edit canonical lessons, then run npm run generate:books.
% canonical-source-hash: fnv1a64:8dfa8697ed1caf20
% canonical-lessons: FR-C418-embarquer, FR-C418-hotesse, FR-C418-steward, FR-C418-hublot, FR-C418-allee

\chapter{À bord}
\label{ch:fr-a-bord}
\hlchaptermodality{418}

\begin{chapteropening}
\textbf{By the end of this chapter:} \emph{I can talk about life on board: boarding, the flight attendants, the window and the aisle.}

Complete the last lesson of chapter 418: I can talk about life on board: boarding, the flight attendants, the window and the aisle.
\end{chapteropening}

\section[embarquer]{embarquer — to board; to take along}

\label{lesson:FR-C418-embarquer}



\begin{quote}
Before the new one: say the French for to search, then the French for to change planes.
\end{quote}

\subsection*{You'll want to know: embarquer}
\textbf{embarquer} — \textquotedblleft{}to board; to take along\textquotedblright{}.

\begin{quote}
Les passagers embarquent par l'avant de l'avion.
\end{quote}

\begin{quote}
\emph{The passengers board at the front of the plane.}
\end{quote}

Informal: \emph{embarquer quelqu'un} — to take someone along.

\begin{cousinweb}
\emph{embarquer} is built from \emph{en-} and \emph{barque} 'small boat', which came through Italian or Provençal from Latin \emph{barca} (Lewis \& Short). EtymDB lists the verb without a source.
\end{cousinweb}

\begin{grammarlens}[title={What you've built}]
One more word for telling, explaining and describing at B1.
\end{grammarlens}

\subsection*{Your turn}
\begin{itemize}
\raggedright
  \item \emph{Say it:} \emph{embarquer}
  \item \emph{Say it:} \emph{embarquer}, once more
  \item \emph{Say it:} a sentence of your own with \emph{embarquer}
  \item \emph{Recall:} say the French for to search, then the French for to change planes, then say \emph{embarquer} again
\end{itemize}

\begin{tcolorbox}[breakable,colback=teal!4,colframe=teal!35!black,title={Before you move on}]
What does embarquer mean? (\textquotedblleft{}To board; to take along\textquotedblright{}.) Say it once more.
\end{tcolorbox}
\section[l'hôtesse]{l'hôtesse — the flight attendant (a woman); the hostess}

\label{lesson:FR-C418-hotesse}



\begin{quote}
Before the new one: say the French for to change planes, then the French for to board.
\end{quote}

\subsection*{You'll want to know: l'hôtesse}
\textbf{l'hôtesse} — \textquotedblleft{}the flight attendant (a woman); the hostess\textquotedblright{}.

\begin{quote}
L'hôtesse sert des boissons pendant le vol.
\end{quote}

\begin{quote}
\emph{The flight attendant serves drinks during the flight.}
\end{quote}

\emph{Une hôtesse de l'air}; today also \emph{un ou une agent de bord}.

\begin{cousinweb}
\emph{hôtesse} is the feminine of \emph{hôte} 'host', which goes back to Latin \emph{hospes} 'host, guest' (Lewis \& Short). EtymDB does not list the feminine.
\end{cousinweb}

\begin{grammarlens}[title={What you've built}]
One more word for telling, explaining and describing at B1.
\end{grammarlens}

\subsection*{Your turn}
\begin{itemize}
\raggedright
  \item \emph{Say it:} \emph{l'hôtesse}
  \item \emph{Say it:} \emph{l'hôtesse}, once more
  \item \emph{Say it:} a sentence of your own with \emph{l'hôtesse}
  \item \emph{Recall:} say the French for to change planes, then the French for to board, then say \emph{l'hôtesse} again
\end{itemize}

\begin{tcolorbox}[breakable,colback=teal!4,colframe=teal!35!black,title={Before you move on}]
What does l'hôtesse mean? (\textquotedblleft{}The flight attendant (a woman); the hostess\textquotedblright{}.) Say it once more.
\end{tcolorbox}
\section[le steward]{le steward — the steward (on a plane)}

\label{lesson:FR-C418-steward}



\begin{quote}
Before the new one: say the French for to board, then the French for the flight attendant.
\end{quote}

\subsection*{You'll want to know: le steward}
\textbf{le steward} — \textquotedblleft{}the steward (on a plane)\textquotedblright{}.

\begin{quote}
Le steward aide une dame avec son sac.
\end{quote}

\begin{quote}
\emph{The steward helps a lady with her bag.}
\end{quote}

Pronounced more or less as in English.

\begin{cousinweb}
\emph{steward} was borrowed from English \emph{steward} (EtymDB).
\end{cousinweb}

\begin{grammarlens}[title={What you've built}]
One more word for telling, explaining and describing at B1.
\end{grammarlens}

\subsection*{Your turn}
\begin{itemize}
\raggedright
  \item \emph{Say it:} \emph{le steward}
  \item \emph{Say it:} \emph{le steward}, once more
  \item \emph{Say it:} a sentence of your own with \emph{le steward}
  \item \emph{Recall:} say the French for to board, then the French for the flight attendant, then say \emph{le steward} again
\end{itemize}

\begin{tcolorbox}[breakable,colback=teal!4,colframe=teal!35!black,title={Before you move on}]
What does le steward mean? (\textquotedblleft{}The steward (on a plane)\textquotedblright{}.) Say it once more.
\end{tcolorbox}
\section[le hublot]{le hublot — the porthole, the window (of a plane or ship)}

\label{lesson:FR-C418-hublot}



\begin{quote}
Before the new one: say the French for the flight attendant, then the French for the steward.
\end{quote}

\subsection*{You'll want to know: le hublot}
\textbf{le hublot} — \textquotedblleft{}the porthole, the window (of a plane or ship)\textquotedblright{}.

\begin{quote}
Ma fille veut toujours une place côté hublot.
\end{quote}

\begin{quote}
\emph{My daughter always wants a window seat.}
\end{quote}

\emph{Côté hublot} — window seat.

\begin{cousinweb}
The origin of \emph{hublot} is not settled, and EtymDB does not list it. It was first a sailors' word.
\end{cousinweb}

\begin{grammarlens}[title={What you've built}]
One more word for telling, explaining and describing at B1.
\end{grammarlens}

\subsection*{Your turn}
\begin{itemize}
\raggedright
  \item \emph{Say it:} \emph{le hublot}
  \item \emph{Say it:} \emph{le hublot}, once more
  \item \emph{Say it:} a sentence of your own with \emph{le hublot}
  \item \emph{Recall:} say the French for the flight attendant, then the French for the steward, then say \emph{le hublot} again
\end{itemize}

\begin{tcolorbox}[breakable,colback=teal!4,colframe=teal!35!black,title={Before you move on}]
What does le hublot mean? (\textquotedblleft{}The porthole, the window (of a plane or ship)\textquotedblright{}.) Say it once more.
\end{tcolorbox}
\section[l'allée]{l'allée — the aisle; the path, the drive}

\label{lesson:FR-C418-allee}



\begin{quote}
Before the new one: say the French for the steward, then the French for the porthole.
\end{quote}

\subsection*{You'll want to know: l'allée}
\textbf{l'allée} — \textquotedblleft{}the aisle; the path, the drive\textquotedblright{}.

\begin{quote}
Il préfère une place côté allée.
\end{quote}

\begin{quote}
\emph{He prefers an aisle seat.}
\end{quote}

A feminine noun. In a park, \emph{une allée} is a path lined with trees.

\begin{cousinweb}
\emph{allée} is built on \emph{aller} 'to go' (EtymDB): a place where you walk.
\end{cousinweb}

\begin{grammarlens}[title={What you've built}]
One more word for telling, explaining and describing at B1.
\end{grammarlens}

\subsection*{Your turn}
\begin{itemize}
\raggedright
  \item \emph{Say it:} \emph{l'allée}
  \item \emph{Say it:} \emph{l'allée}, once more
  \item \emph{Say it:} a sentence of your own with \emph{l'allée}
  \item \emph{Recall:} say the French for the steward, then the French for the porthole, then say \emph{l'allée} again
\end{itemize}

\begin{tcolorbox}[breakable,colback=teal!4,colframe=teal!35!black,title={Before you move on}]
What does l'allée mean? (\textquotedblleft{}The aisle; the path, the drive\textquotedblright{}.) Say it once more.
\end{tcolorbox}
